\subsection{Admitting the opening that the selected program will use}
\label{sec:declared-opening-binding}

A balanced, typed opening book can still violate the selected company's
operating conventions. Its claim map may be empty while receivables are
nonzero; its unit-cost leaf may have the wrong dimensions; its account roles
may overlap; or its statement mapping may omit an account. The opening check
must establish these additional conditions before events run. Yet it must not
silently promise that every future command is valid: an unknown future customer
or an unsupported event is still decided at its execution boundary.

\code{DeclaredOpeningAdmission/1} records the retained empty-opening contract
as ordered declarations. Its template specifies the input envelope, policy
fields and values, exact claim-row and clock types, empty claim population,
cursor equal to opening date, complete and distinct account roles, required
account classes and cash identity. It then checks the unit-cost pointer,
dimensions and strictly positive value, the opening receivable total, the
opening statement report and the command-list limit. Each condition retains
its native position and failure reason. All template stages compile before
their evaluation; a malformed later declaration cannot disappear because an
earlier condition would refuse the opening.

Structural conditions compare declared record fields, schema nodes, role maps
and literal policy values. Scalar date and monetary predicates use the existing
guarded requirements evaluator, retaining their dimensional types until that
evaluator lowers them. This separation preserves a subtle input boundary:
a malformed carry-policy value fails the declared policy equality, rather than
acquiring an earlier and different scalar-type error. Types remain essential
where the selected contract requires typed numerical or date comparisons.

The empty-claim premise makes the opening claim-total condition particularly
simple: the receivables balance must be zero. This is not a replacement for
the row invariants of an operating book. Once events create claims, the complete
authored state invariant again checks their totals, dates, references and
assessment snapshots. The opening report retains its own evaluator, which
checks the period start, opening flows and temporary balances, statement-map
coverage and statement arithmetic. No statement rule is duplicated inside the
opening template.

Two templates reproduce exact acceptance or refusal on 82 complete inputs,
including the 58 preceding audit cases. Ten are admitted. Among them are
inputs with a malformed future command, unknown operation or unknown customer
that the original opening validator also accepts. The corresponding execution
checks remain separate. A deliberate change from strictly positive to
nonnegative unit cost admits a zero-cost opening; this proves that the declared
predicate controls the decision and explicitly differs from the selected native
contract. Dependency identities, input order and stale selections are checked
alongside the substantive conditions.

Passing opening admission and passing execution independently would still
leave a serious gap. One could admit a book under one account-role map but
execute methods selected for another, or admit a different command list from
the one that actually runs. \code{AdmittedDeclarativeProgram/1} therefore binds
both evaluations to one explicit selection. That selection names the chart
context, role map, unit-cost pointer, policy, opening template, complete rule set
and state invariant. The call binds all supplied definitions and facts rather
than accepting a caller's success flag.

After admission, the adapter compares the complete shared opening, schemas,
statement mapping and memory definition. It checks the selected rules and
invariant, then compares every projected event, report and carry command.
Projection removes only the assessment-key order and memory carry-policy fields
that the native-shaped admission input lacks; the composer checks those fields
when they are used. Event-input order is retained explicitly because it affects
the ordered list of unknown inputs. A mismatch refuses execution.

With all native lifecycle functions and the operating-method builder disabled,
the combined path preserves the complete prior composer results for four
successful programs and 26 adverse programs. The successful programs contain
44 commands whose complete substantive payloads already agree with the native
baseline. Twenty-two binding controls reject changed facts, roles, cost
selection, policy, rules, invariant, commands, order and stale identities.
An opening refusal executes no command. Existing typed-memory, transition,
requirements, report and carry evaluators remain the owners of their respective
calculations.

This removes the native opening evaluator prerequisite for the selected
combined path. The selection is still a finite reviewed chart contract, not
a proof that arbitrary edited rules have equivalent meanings. The portable
inputs use exact JSON values and rational strings. Earlier company-source
frontends retain their older target. Section~\ref{sec:selected-company-source}
introduces the successor that expresses these algorithms directly. General nonempty opening populations, wider company rules,
independent implementation and accounting acceptance remain open.

\begingroup\raggedright
The opening comparisons and bound programs are retained under
\path{data/accounting-master/declared-opening-admission-20261026/}
and \path{data/accounting-master/admitted-program-binding-20261026/}.
\par
\endgroup
